% Options for packages loaded elsewhere
\PassOptionsToPackage{unicode}{hyperref}
\PassOptionsToPackage{hyphens}{url}
\documentclass[
  12pt,
]{ctexart}
\usepackage{xcolor}
\usepackage{amsmath,amssymb}
\setcounter{secnumdepth}{-\maxdimen} % remove section numbering
\usepackage{iftex}
\ifPDFTeX
  \usepackage[T1]{fontenc}
  \usepackage[utf8]{inputenc}
  \usepackage{textcomp} % provide euro and other symbols
\else % if luatex or xetex
  \usepackage{unicode-math} % this also loads fontspec
  \defaultfontfeatures{Scale=MatchLowercase}
  \defaultfontfeatures[\rmfamily]{Ligatures=TeX,Scale=1}
\fi
\usepackage{lmodern}
\ifPDFTeX\else
  % xetex/luatex font selection
\fi
% Use upquote if available, for straight quotes in verbatim environments
\IfFileExists{upquote.sty}{\usepackage{upquote}}{}
\IfFileExists{microtype.sty}{% use microtype if available
  \usepackage[]{microtype}
  \UseMicrotypeSet[protrusion]{basicmath} % disable protrusion for tt fonts
}{}
\makeatletter
\@ifundefined{KOMAClassName}{% if non-KOMA class
  \IfFileExists{parskip.sty}{%
    \usepackage{parskip}
  }{% else
    \setlength{\parindent}{0pt}
    \setlength{\parskip}{6pt plus 2pt minus 1pt}}
}{% if KOMA class
  \KOMAoptions{parskip=half}}
\makeatother
\usepackage{longtable,booktabs,array}
\usepackage{calc} % for calculating minipage widths
% Correct order of tables after \paragraph or \subparagraph
\usepackage{etoolbox}
\makeatletter
\patchcmd\longtable{\par}{\if@noskipsec\mbox{}\fi\par}{}{}
\makeatother
% Allow footnotes in longtable head/foot
\IfFileExists{footnotehyper.sty}{\usepackage{footnotehyper}}{\usepackage{footnote}}
\makesavenoteenv{longtable}
\setlength{\emergencystretch}{3em} % prevent overfull lines
\providecommand{\tightlist}{%
  \setlength{\itemsep}{0pt}\setlength{\parskip}{0pt}}
\usepackage{bookmark}
\IfFileExists{xurl.sty}{\usepackage{xurl}}{} % add URL line breaks if available
\urlstyle{same}
\hypersetup{
  hidelinks,
  pdfcreator={LaTeX via pandoc}}

\author{SCX}
\date{2026}

\begin{document}

\section{A2
的严格数学分析}\label{a2-ux7684ux4e25ux683cux6570ux5b66ux5206ux6790}

\begin{quote}
触发: hostile review 指出 \(\exp(-2M\Delta^2/(1+(M-1)\rho))\)
不是严格定理 结论: 正确。方差膨胀代入 Hoeffding
是指启发式，不是严格界。但 A2 本身有更强的论证。
\end{quote}

\begin{center}\rule{0.5\linewidth}{0.5pt}\end{center}

\subsection{\texorpdfstring{1. \(\exp(-2M t^2/(1+(M-1)\rho))\)
是不是严格定理？}{1. \textbackslash exp(-2M t\^{}2/(1+(M-1)\textbackslash rho)) 是不是严格定理？}}\label{exp-2m-t21m-1rho-ux662fux4e0dux662fux4e25ux683cux5b9aux7406}

\textbf{不是。}

Hoeffding 不等式的证明依赖:
\[\mathbb{E}\left[\exp\left(\lambda\sum_{i=1}^M (X_i - \mu_i)\right)\right] = \prod_{i=1}^M \mathbb{E}[\exp(\lambda(X_i - \mu_i))]\]

这个\textbf{因式分解}要求 \(\{X_i\}\) 独立。如果 \(X_i\) 相关，MGF
不能因式分解，你不能简单地把 \(M\) 替换成 \(M/(1+(M-1)\rho)\)
就得到严格指数界。

方差膨胀 \((1+(M-1)\rho)\) 对 \textbf{\(\text{Var}(\sum X_i)\)}
是正确的，但对 \textbf{指数集中} 不正确。方差只给 Chebyshev
界（多项式速率 \(1/t^2\)），不是指数界。

\textbf{结论}: 方差膨胀 Hoeffding
是启发式（heuristic），在调查抽样和群组随机试验中广泛使用，但不是数学定理。把它写进
SI 会被审稿人抓。

\begin{center}\rule{0.5\linewidth}{0.5pt}\end{center}

\subsection{2. 那 A2
本身到底站不站得住？}\label{ux90a3-a2-ux672cux8eabux5230ux5e95ux7ad9ux4e0dux7ad9ux5f97ux4f4f}

\subsubsection{2.1 更强的论证: A1 蕴含
A2}\label{ux66f4ux5f3aux7684ux8bbaux8bc1-a1-ux8574ux542b-a2}

A1 说: \(M\) 个专家在\textbf{不相交的独立同分布子集}上训练。
\[D_1, \dots, D_M \sim_{\text{i.i.d.}} P^n \quad \text{且} \quad D_i \cap D_j = \varnothing\]

每个 \(f_m = \mathcal{A}(D_m)\) 是训练算法 \(\mathcal{A}\)
作用于数据子集 \(D_m\) 的\textbf{确定函数}。

因为 \(D_1, \dots, D_M\) 是独立随机变量，\(f_1, \dots, f_M\)
也是独立的随机函数。

对任意固定的测试点 \(x\):
\[e_m(x, y) = \mathbf{1}\{\ell(f_m(x), y) > \tau\}\]

因为 \(f_m\) 只依赖 \(D_m\)，且
\(D_1 \perp D_2 \perp \cdots \perp D_M\):
\[e_1 \perp e_2 \perp \cdots \perp e_M \quad \mid \quad x\]

\textbf{A2 不是假设------它是 A1 的推论。}

\subsubsection{2.2
概率空间是什么？}\label{ux6982ux7387ux7a7aux95f4ux662fux4ec0ux4e48}

Hoeffding
界的概率空间是\textbf{训练数据的随机性}，不是测试数据的随机性。

\[\mathbb{P}_{D_1,\dots,D_M}\left(C(x) > \theta \mid \text{clean}, x\right) \leq \exp(-2M(\theta-\mu_s)^2)\]

意思是: 如果你随机划分训练数据 \(M\) 次、独立训练 \(M\)
个专家，得到的专家集以高概率具有好的噪声检测性质。

\textbf{训练完成后}，你有一组固定的专家。对于给定的测试样本
\(x\)，\(C(x)\) 是一个确定性的数。Hoeffding 不适用。

但这不削弱定理------所有的泛化界都有这个性质。定理告诉你''这个\textbf{方法}是好的''，不是''这组\textbf{特定专家}是好的''。

\subsubsection{2.3 Hostile reviewer
的反驳为什么不对}\label{hostile-reviewer-ux7684ux53cdux9a73ux4e3aux4ec0ux4e48ux4e0dux5bf9}

\begin{quote}
``所有 CNN 在模糊图像上都一起错''
\end{quote}

这是在说: \textbf{给定已经训练好的
CNN}，它们在模糊图像上的确定行为是一致的。

但 Hoeffding 问的是:
\textbf{在训练之前}，随机划分数据并独立训练后，你有多大把握得到一组能检测噪声的专家？

两者不在同一个概率空间里。Hostile reviewer
混淆了\textbf{训练后确定性行为}和\textbf{训练前概率保证}。

\subsubsection{2.4
但仍有一个残留问题}\label{ux4f46ux4ecdux6709ux4e00ux4e2aux6b8bux7559ux95eeux9898}

即使 A2 严格成立（作为 A1 的推论），\textbf{估计} \(\mu_s\)
仍然需要清洁标签。在有限样本下，\(\hat{\mu}_s\) 的误差会转化为
\(\hat{\Delta}_s\)
的误差。这是\textbf{有限样本估计问题}，不是\textbf{独立性假设问题}。推论
4（有限样本校正）已经用保守上界处理了这个问题。

\begin{center}\rule{0.5\linewidth}{0.5pt}\end{center}

\subsection{3.
诚实的修改建议}\label{ux8bdaux5b9eux7684ux4feeux6539ux5efaux8bae}

\subsubsection{3.1 不要引入 A2'}\label{ux4e0dux8981ux5f15ux5165-a2}

方差膨胀 Hoeffding 是启发式，不够格写进定理陈述。保持原有的
A2，但\textbf{在正文和 SI 中诚实讨论}。

\subsubsection{3.2 在 SI 中加一节 ``On the Role of Assumption
A2''}\label{ux5728-si-ux4e2dux52a0ux4e00ux8282-on-the-role-of-assumption-a2}

内容: 1. A2 实际上是 A1 的推论（不相交训练集 → 独立专家 → 独立错误） 2.
概率空间是训练随机性，不是测试行为 3. 有限样本下 \(\mu_s\)
的估计误差有独立处理（推论 4） 4.
在实践中，专家错误可能表现出表面相关（同一张模糊图），但这反映的是特征相似性导致相似的
\(\mu_s\)，而非统计相关性 5. 对 A2 的经验诊断:
在验证集上估计专家错误的条件相关系数；如果显著偏离零，检查训练集是否真的不相交

\subsubsection{3.3
在正文中诚实措辞}\label{ux5728ux6b63ux6587ux4e2dux8bdaux5b9eux63aaux8f9e}

把 ``conditional independence'' 改为: \textgreater{} ``conditional
independence of expert errors, which follows from training on disjoint
data subsets (Assumption A1)''

这样把 A2 从''假设''降级为''A1 的逻辑结果''，同时指出它依赖于 A1
的正确执行。

\begin{center}\rule{0.5\linewidth}{0.5pt}\end{center}

\subsection{4. 最终判断}\label{ux6700ux7ec8ux5224ux65ad}

\begin{longtable}[]{@{}
  >{\raggedright\arraybackslash}p{(\linewidth - 2\tabcolsep) * \real{0.6000}}
  >{\centering\arraybackslash}p{(\linewidth - 2\tabcolsep) * \real{0.4000}}@{}}
\toprule\noalign{}
\begin{minipage}[b]{\linewidth}\raggedright
问题
\end{minipage} & \begin{minipage}[b]{\linewidth}\centering
状态
\end{minipage} \\
\midrule\noalign{}
\endhead
\bottomrule\noalign{}
\endlastfoot
A2 是否站得住？ & \textbf{是} --- A1 蕴含 A2（训练集不相交 →
专家独立） \\
在实践中是否违反？ & \textbf{取决于 A1 的执行} ---
如果训练集真的不相交，A2 成立 \\
\(\exp(-2M\Delta^2/(1+(M-1)\rho))\) 是否严格？ & \textbf{否} ---
这是启发式，不应写进定理 \\
应该如何修改？ & 澄清 A2 是 A1 的推论；在 SI
中加入诚实讨论；不要引入虚假的 A2' \\
对论文的影响程度 & \textbf{低} ---
不需要改定理陈述，只需要加一段诚实讨论 \\
\end{longtable}

\begin{center}\rule{0.5\linewidth}{0.5pt}\end{center}

\subsection{5. 对 hostile review
其他批评的回应}\label{ux5bf9-hostile-review-ux5176ux4ed6ux6279ux8bc4ux7684ux56deux5e94}

\subsubsection{``CIFAR F1=0.617, 理论下界
F1≥0.976''}\label{cifar-f10.617-ux7406ux8bbaux4e0bux754c-f10.976}

这已经在 §4.3（紧致性讨论）中诚实处理了: \(\mu_s\) 在 CIFAR 实验中约为
0.45（3-epoch CPU 训练后），而非假定的 0.20。代入正确的 \(\mu_s\)
后下界为 F1≥0.18，与经验值 0.617 一致。

\subsubsection{\texorpdfstring{``\(\mu_s\) 需要清洁标签 →
鸡生蛋''}{``\textbackslash mu\_s 需要清洁标签 → 鸡生蛋''}}\label{mu_s-ux9700ux8981ux6e05ux6d01ux6807ux7b7e-ux9e21ux751fux86cb}

冷启动协议需要少量锚点样本（推论 1），这是理论上的必要成本（Thm 3
证明了没有锚点时不可识别）。这不是循环定义------是对不可识别性定理的工程回应。

\subsubsection{``Bootstrap 诊断的 τ=0.7
是拍脑袋''}\label{bootstrap-ux8bcaux65adux7684-ux3c40.7-ux662fux62cdux8111ux888b}

承认。\(\tau=0.7\) 是初始校准值，应在具体领域中根据已知噪声标签校准。SI
中应明确标注为''建议初始值，需领域校准''。

\begin{center}\rule{0.5\linewidth}{0.5pt}\end{center}

\emph{关联: {[}{[}01\_noise\_detection\_guarantee{]}{]} ·
{[}{[}03\_unidentifiability{]}{]} ·
{[}{[}a2\_correlation\_analysis{]}{]}}

\end{document}
